% !TEX program = lualatex
\documentclass[12pt, a4paper]{article}
\usepackage{master}

\DeclareWorkGR{Cat}{범주론}{Κατηγορίαι}{Categoriae}
\DeclareWorkGR{DI}{명제론}{Περὶ ἑρμηνείας}{De interpretatione}
\DeclareWorkGR{APr}{분석론 전서}{Ἀναλυτικὰ πρότερα}{Analytica priora}
\DeclareWorkGR{DA}{영혼론}{Περὶ ψυχῆς}{De anima}

\fancyhead[L]{\footnotesize\color{midgray}오르가논 강독}
\fancyhead[R]{\footnotesize\color{midgray}특별 세션 γ}

\begin{document}
\thispagestyle{firstpage}

\weeksection{특별 세션 γ}{명제론 종합}{의미에서 양상까지, 세 저작의 합류}

\subsection{세 저작의 논리적 상승}

\subsubsection{비결합에서 결합으로, 결합에서 양상으로}

범주론(2--5주차)에서 우리는 비결합 표현들을 분석했습니다. 이사고게(6--8주차)에서 이 단순 항들이 주어와 술어로서 관계하는 다섯 가지 방식을 분석했습니다. 명제론(9--11주차)에서 이 항들이 실제로 결합되어 명제를 이루었을 때 참이거나 거짓이 되는 논리적 구조를 탐구했습니다. 이 상승은 네 단계로 정리됩니다: (1) 비결합 표현(범주론, 10범주) → (2) 술어화 방식(이사고게, 다섯 술어가능자) → (3) 명제(명제론, 대당 관계·양상) → (4) 추론(분석론 전서, 삼단논법).

\subsubsection{범주론의 두 축이 명제론으로 흘러들어오는 길}

범주론의 두 기본 축 --- `주어에 대해 말해짐'(본질적 술어화)과 `주어 안에 있음'(부수적 내재) --- 은 명제론의 구조를 결정합니다. 첫째, 긍정과 부정(6장)은 무엇인가를 무엇인가에 대해 `결합하거나 분리하는 것'입니다 --- 이것은 범주론 4장(1a16--19)의 `결합/비결합' 구별을 명제 차원으로 격상시킨 것입니다. 둘째, 보편자와 단칭자의 구별(7장)은 범주론의 제1실체/제2실체 구별에 대응합니다. 셋째, 반대자 수용(3주차, 범주론 5장)은 해전 논변(9장)의 배경에 있습니다: 실체가 반대자를 수용할 수 있다는 것은 명제의 진리값이 시간에 따라 변할 수 있다는 것을 함축합니다. 넷째, 후범주론(5주차)의 네 가지 대립은 대당 사각형의 반대·모순으로 변환됩니다.

\subsubsection{이사고게의 도구가 명제론에서 작동하는 방식}

이사고게의 다섯 술어가능자는 명제의 양상적 지위를 결정합니다 --- 비록 아리스토텔레스 자신이 이것을 명시적으로 연결하지는 않지만, 그 연결은 분명합니다. 류·종차·고유성 술어는 필연적 참인 A-명제를 만들고(`모든 사람은 동물이다', `모든 사람은 이성적이다', `모든 사람은 웃을 수 있다'), 분리 가능한 부수적 술어는 우연적인 A-명제를 만듭니다(`모든 사람은 앉아 있다'). 술어의 종류가 명제의 양상적 지위를 결정한다는 것은 12--13장의 양상 명제 이론과 분석론 후서의 `그 자체로'(\gr{καθ' αὑτό}) 술어화 이론을 예비합니다.

{\footnotesize
\begin{concepttable}{|l|l|c|c|}
\hline
\rowcolor{tableheadblue}
{\color{h1navy}\bfseries 술어가능자} &
{\color{h1navy}\bfseries 예시 A-명제} &
{\color{h1navy}\bfseries 진리값} &
{\color{h1navy}\bfseries 양상} \\ \hline
류 & 모든 사람은 동물이다 & 필연적 참 & 필연 \\ \hline
종차 & 모든 사람은 이성적이다 & 필연적 참 & 필연 \\ \hline
고유성 & 모든 사람은 웃을 수 있다 & 필연적 참 & 필연 \\ \hline
부수성 (분리불가) & 모든 까마귀는 검다 & 우연적 참 & 우연 \\ \hline
부수성 (분리가능) & 모든 사람은 앉아 있다 & 우연적 거짓 & 우연 \\ \hline
\end{concepttable}
}

\subsection{명제론의 핵심 성취와 열린 문제들}

\subsubsection{의미의 삼각형 --- 네 가지 미해결 물음}

명제론 1장(16a3--8)의 의미의 삼각형은 서양 의미론의 출발점이지만, 네 가지 미해결 물음을 남깁니다. 첫째, 영혼의 경험(\gr{παθήματα τῆς ψυχῆς})의 정체 --- 모드락의 심상 해석, 애크릴의 회의적 독해, 크레츠만의 `이론 없음' 독해 사이의 논쟁은 해결되지 않았습니다. 둘째, 관습성의 범위 --- 문장의 의미론적 구조(주어-술어-계사)도 관습적인가, 자연적 구조를 반영하는 것인가? 셋째, 의미의 삼각형과 심플리키오스의 `세 가지는 원래 하나였다'(특별 세션 α)의 관계. 넷째, `염소사슴'(\gr{τραγέλαφος}, 16a16) --- 존재하지 않는 것을 지시하는 이름이 의미한다는 것은 무엇인가?

\subsubsection{대당 사각형 --- 2500년의 구조물}

대당 사각형은 아리스토텔레스 논리학에서 가장 오래 살아남은 단일 구조물입니다. 2세기 아풀레이우스로부터 시작하여, 보에티우스를 거쳐 중세 대학의 표준 교재가 되었고, 19세기까지 모든 논리학 교과서에 등장했습니다. 파슨스(2025)의 방어가 보여주듯이, 현대 기호 논리학에 의한 비판은 보에티우스의 O-형식 변형에 기인하며, 아리스토텔레스의 원래 체계는 정합적이었습니다.\footnote{SEP, ``The Traditional Square of Opposition'' (Parsons, rev.\ 2025).}

사고 연습: 대당 사각형과 포르피리오스의 나무를 결합해 봅시다. 포르피리오스의 나무가 실체 범주의 내부 구조(류-종-종차 위계)를 보여주듯이, 대당 사각형은 명제의 양과 질의 내부 구조(A-E-I-O)를 보여줍니다. 둘 다 2×2 분류 체계입니다. 이 두 구조가 만나는 지점은 분석론 전서의 삼단논법입니다.

\subsubsection{해전 논변 --- 왜 2500년간 논의되는가}

해전 논변이 가장 오래 논의된 논변 가운데 하나인 이유는, 세 가지 근본적인 철학적 물음을 한꺼번에 건드리기 때문입니다. 첫째, 진리와 시간의 관계 --- 명제의 진리값은 시간에 의존하는가? 둘째, 결정론과 자유 의지 --- 모든 미래 명제가 참이거나 거짓이라면, 미래는 이미 확정되어 있고 숙고는 무의미합니다. 셋째, 논리적 원리의 범위 --- 이가 원리와 배중률은 같은 것인가, 다른 것인가?

{\footnotesize
\begin{concepttable}{|l|l|l|}
\hline
\rowcolor{tableheadblue}
{\color{h1navy}\bfseries 물음} &
{\color{h1navy}\bfseries 아리스토텔레스의 답} &
{\color{h1navy}\bfseries 후속 전통} \\ \hline
진리와 시간 & 진리값은 시간에 따라 변함 & 시제 논리 (Prior 1957) \\ \hline
결정론/자유 의지 & 미래는 열려 있다 (18b31) & 보에티우스, 아퀴나스, 오캄 \\ \hline
논리적 원리 & 배중률 보존, 이가 원리 제한 & 우카시에비치, 맥팔레인 \\ \hline
\end{concepttable}
}

\subsubsection{계사의 발견과 양상 명제의 이중성}

10장에서 `있다'가 계사와 존재 술어로 구별된 것은, `있다는 무엇을 의미하는가?'라는 서양 존재론의 핵심 물음의 출발점입니다. 12--13장의 \gr{δυνατόν}의 이중성 --- 양면적 가능성(우연성)과 단면적 가능성(불가능하지 않음) --- 은 분석론 전서의 양상 삼단논법에 직접 연결됩니다. 양상 동치(□p ↔ ¬◇¬p)\footnote{양상 논리 기호: □ = `필연적으로', ◇ = `가능하게', ¬ = `아닌', ↔ = `서로 같다'. □p ↔ ¬◇¬p는 `p가 필연적이라는 것은, p가 아닌 것이 가능하지 않다는 것과 같다'로 읽습니다.}는 현대 양상 논리의 쌍대성 법칙을 2천 년 앞서 진술한 것입니다.

\subsection{고대 주석 전통에서의 명제론}

\subsubsection{암모니오스의 명제론 주석}

암모니오스(c.\ 435--517/526)의 명제론 주석은 특별 세션 α에서 다룬 범주론 주석과 함께 그의 주요 저작 가운데 하나입니다.\footnote{Ammonius, \gri{In Aristotelis De interpretatione commentarius}, ed.\ A.\ Busse, CAG IV.5 (Berlin, 1897).} 16a3--8의 의미의 삼각형에 대해 `진리의 담지자'가 무엇인지를 물었습니다 --- 사물의 실재적 결합/분리인가, 영혼의 경험인가, 발화된 진술인가? 이것은 범주론의 \gr{σκοπός} 논쟁의 명제론적 변형입니다.

\subsubsection{보에티우스의 이중 주석}

보에티우스는 명제론에 대해서도 이사고게와 마찬가지로 두 편의 주석서를 남겼습니다.\footnote{John Magee, \gri{Boethius on Signification and Mind} (Leiden: Brill, 1989).} 핵심 공헌: 대당 사각형의 도식화(아풀레이우스의 도식을 전달·정교화), 9장의 미래 우연 명제에 대한 `비확정적 참' 해석, 계사와 존재 술어의 구별에 대한 라틴어 용어 확립. 특별 세션 β에서 다루었듯이, 보에티우스의 명제론 주석은 구논리학의 핵심 구성 요소였습니다.

\subsubsection{알파라비의 해전 논변 독해}

알파라비의 명제론 주석에서 해전 논변에 대한 입장은 학술적으로 논쟁적입니다. 레셔(Rescher)는 알파라비가 미래 우연 명제의 진리값이 이미 `비확정적 방식으로' 분배되어 있다고 해석하고(비표준적 해석에 가까움), 짐머만(Zimmerman)은 진리값이 아직 분배되지 않았다고 해석합니다(전통적 해석에 가까움). 특별 세션 β에서 다룬 알파라비의 1차/2차 개념 구별을 상기합시다: 명제론의 긍정·부정·모순 같은 개념은 전형적인 2차 개념이므로, 명제론은 범주론보다 더 확실하게 논리학에 속합니다.

\subsection{명제론 읽기에서 우리가 배운 것}

\subsubsection{명제론이 변증술을 위한 준비라는 것}

휘태커(1996)의 핵심 논제를 최종적으로 정리합시다: 명제론은 분석론을 위한 예비가 아니라, 토피카와 소피스트적 논박(변증술)을 위한 예비입니다. 근거 세 가지: (1) 중심 주제가 모순 대립이며, 모순 대립의 확립은 변증술적 논박의 핵심 과제; (2) 8장의 동음이의 분석이 소피스트적 논박과 직접 병행; (3) 14장의 결론(반대 믿음은 모순적 진술로 표현)이 변증술적 실천의 기초. 그러나 이것이 분석론과 무관하다는 뜻은 아닙니다 --- 대당 관계(특히 환위)는 \citework{APr}{I.2, 25a1--25}에서 삼단논법 이론의 직접적인 전제로 사용됩니다. 명제론은 변증술과 분석술 모두를 위한 공통 기초입니다.

\subsubsection{분석론 전서로의 전환}

다음 주부터 시작하는 분석론 전서는 명제론의 명제들이 전제(\gr{πρότασις})로서 삼단논법(\gr{συλλογισμός}) 안에 배치될 때 어떤 타당한 추론이 가능한지를 탐구합니다. \citework{APr}{I.1, 24a16--17}에서 아리스토텔레스는 삼단논법을 `어떤 것들이 놓이면 그것들과 다른 무엇인가가 그 놓인 것들에 의해 필연적으로 따라나오는 논증'이라고 정의합니다.

여기서 마지막 사고 연습을 해봅시다. 지금까지 우리가 읽은 세 저작에서 분석론 전서에 직접 필요한 것들을 추출해 보세요: 범주론에서 --- 10범주(삼단논법의 술어 내용), `주어에 대해 말해짐'(보편적 술어화의 기초); 이사고게에서 --- 류-종-종차 위계(정의와 논증의 구조), 다섯 술어가능자(토피카의 논변 유형); 명제론에서 --- A/E/I/O 명제(삼단논법의 전제 형식), 대당 관계(환위 규칙), 양상 연산자(양상 삼단논법의 전제). 우리가 읽은 모든 것이 분석론 전서를 향해 수렴합니다.

\subsection*{참고문헌}
\begin{references}

\refitem Ackrill, J.\ L., trans. \gri{Aristotle's Categories and De Interpretatione}. Oxford: Clarendon Press, 1963.

\refitem Crivelli, Paolo. \gri{Aristotle on Truth}. Cambridge: Cambridge University Press, 2004.

\refitem Kretzmann, Norman. ``Aristotle on Spoken Sound Significant by Convention.'' In \gri{Ancient Logic and Its Modern Interpretations}, edited by J.\ Corcoran, 3--21. Dordrecht: Reidel, 1974.

\refitem Magee, John. \gri{Boethius on Signification and Mind}. Leiden: Brill, 1989.

\refitem Modrak, Deborah K.\ W. \gri{Aristotle's Theory of Language and Meaning}. Cambridge: Cambridge University Press, 2001.

\refitem Sedley, David. ``Aristotle's \gri{De interpretatione} and Ancient Semantics.'' In \gri{Knowledge Through Signs}, edited by G.\ Manetti, 87--108. Turnhout: Brepols, 1996.

\refitem Whitaker, C.\ W.\ A. \gri{Aristotle's De Interpretatione: Contradiction and Dialectic}. Oxford: Clarendon Press, 1996.

\end{references}

\end{document}
